% paper2_sandpiles_and_lattice_entropy.tex
% REVTeX 4-2, Physical Review E style.
% Scope: Formulas 1, 25-28, 30-31, 33-36 of formulas.txt (Formula 2 deliberately excluded).
\documentclass[aps,pre,10pt,twocolumn,superscriptaddress,nofootinbib]{revtex4-2}

\usepackage{amsmath}
\usepackage{amssymb}
\usepackage{booktabs}
\usepackage[colorlinks=true,allcolors=blue]{hyperref}

\newcommand{\E}{\mathbb{E}}
\newcommand{\Cl}{\operatorname{Cl}_2}
\newcommand{\sq}{\sqrt{3}}
\newcommand{\arccosh}{\operatorname{arccosh}}
\newcommand{\arcsinh}{\operatorname{arcsinh}}

\begin{document}

\title{Exact height-one probabilities of three-dimensional Abelian sandpiles,\\
spanning-tree entropies of the snub-square and Cairo lattices,\\
and elementary resistor-network reductions}

\author{Jacob Goodchild}
\email{jacob.l.goodchild@gmail.com}

\date{September 28, 2026}

\begin{abstract}
We compute exact closed forms for three related lattice quantities: the probability $P_1$
that a site of the Abelian sandpile has minimal height, the spanning-tree entropy, and
two-point resistances. All three are controlled by the same Kirchhoff transfer
currents. In three dimensions we obtain $P_1$ for the simple-cubic, body-centred cubic,
face-centred cubic and pyrochlore lattices as polynomials in the corresponding Watson
integral $W$ and in $1/(\pi^2 W)$. For example, on the simple cubic lattice
$P_1=\tfrac43\alpha^3(2-3\alpha)^2=0.054582470221\ldots$ with
$\alpha=(7W+54/(\pi^2W))/36$. For bipartite lattices of girth at least 6 with symmetric
second neighbours we prove $P_1=(z-2)^{z-1}/[z(z-1)^{z-1}]$, which gives $1/12$ for the
honeycomb (known) and the Laves graph, and $2/27$ for diamond. In two dimensions we give exact
$P_1$ for the king, 4.8.8, Union Jack, star, dice, Lieb, snub-square, checkerboard, triakis
and next-nearest-neighbour honeycomb lattices. The spanning-tree entropy of the snub-square
(Shastry--Sutherland) lattice is $z=\tfrac{10}{3\pi}G+\tfrac13\ln(2+\sq)$, and that of its dual,
the Cairo pentagonal lattice, is $\tfrac{20}{9\pi}G+\tfrac29\ln(2+\sq)$, where $G$ is Catalan's
constant. The Shastry--Sutherland network with arbitrary dimer conductance has elementary
resistances and a Clausen-function entropy. The king lattice has
$z=20G/(3\pi)$. We give a single elementary formula for the axis-bond resistance of any
square lattice with diagonal bonds, and we reduce next-nearest-neighbour networks on bipartite
girth-$\ge6$ lattices to one Green function (elliptic for the honeycomb, Joyce's FCC form for
diamond). Every result was identified by PSLQ at 25--60 digits and confirmed independently by
exact finite-volume computations or Monte Carlo simulation. An appendix gives a Hurwitz-zeta
closed form for the added mass of a sphere touching the inside of a spherical container.
\end{abstract}

\maketitle

%=====================================================================
\section{Introduction}
\label{sec:intro}
%=====================================================================

The Abelian sandpile model~\cite{BTW1987,Dhar1990} is the standard exactly solvable model of
self-organized criticality. Majumdar and Dhar showed that the probability $P_1$ that a site
has the minimal height in the stationary state is a finite determinant of lattice Green
functions~\cite{MajumdarDhar1991}. On the square lattice this gives
$P_1=2/\pi^2-4/\pi^3$. Priezzhev computed the higher heights~\cite{Priezzhev1994}. Exact values
have since been obtained on the honeycomb~\cite{AzimiTafreshi2010}, triangular and
hexagonal~\cite{PonceletRuelle2017}, and, very recently, kagome~\cite{Liu2026} lattices. The
Majumdar--Dhar method applies in any dimension. Nevertheless, to our knowledge no closed form for
$P_1$ on a three-dimensional lattice has appeared. The obstacle is the need for exact
resistances beyond nearest neighbours.

We remove that obstacle by identifying every needed resistance in the natural basis
$\{1,W,1/(\pi^2W)\}$ of the relevant Watson integral~\cite{Watson1939,GlasserZucker1977}. The
same machinery, based on Kirchhoff's transfer currents~\cite{Kirchhoff1847} and their spanning-tree
interpretation~\cite{BurtonPemantle1993,LyonsPeres2016}, also gives new spanning-tree entropies.
Numerical values for many Archimedean lattices are known from Chang and Wang~\cite{ChangWang2006},
Chang and Shrock~\cite{ChangShrock2006}, and Guttmann and Rogers~\cite{GuttmannRogers2012}; we add
closed forms. It also gives exact resistances on networks whose resistances were previously known
only numerically~\cite{Owaidat2013} or via mappings~\cite{Cserti2011}.

\emph{Status of novelty.} All statements were checked by web search against the literature. The
honeycomb and kagome values of $P_1$ are known~\cite{AzimiTafreshi2010,Liu2026}; we reproduce them
as validation. The king-lattice entropy $20G/(3\pi)$ may be implicit in the Mahler-measure
literature~\cite{GuttmannRogers2012}. For the remaining results our searches found no prior closed
form. Novelty should be read in that limited sense.

\subsection*{Notation}
All resistors are $1\,\Omega$ unless a conductance is specified, and $R(x)$ is the effective
resistance between the origin and $x$. For an arc $a=(o\to h_a)$ the transfer current is
$Y_{ab}=\tfrac12[R(h_a)+R(h_b)-R(h_a-h_b)]$. $G$ denotes Catalan's constant, and
$\Cl(\theta)=\sum_{n\ge1}\sin(n\theta)/n^2$ is the Clausen function. Complete elliptic integrals
use the parameter convention $K(m)=\int_0^{\pi/2}(1-m\sin^2\theta)^{-1/2}\,d\theta$.
$\E_k$ denotes the normalized Brillouin-zone average.

%=====================================================================
\section{Abelian sandpiles on 3D lattices}
\label{sec:3d}
%=====================================================================

\subsection{Method}

In the Abelian sandpile on a graph, a site $v$ topples when its height exceeds $\deg(v)$ and
sends one grain to each neighbour. Majumdar and Dhar~\cite{MajumdarDhar1991} showed that
\begin{equation}
P_1=\det(I-M),\qquad M_{jk}=D(x_j)+D(x_k)-D(x_j-x_k),
\label{eq:md}
\end{equation}
where $D=R/2$ and $x_j$ runs over the $z-1$ neighbours of the origin whose bonds are removed
by the defect (one bond is kept). The matrix $M$ is exactly the transfer-current matrix $Y$
restricted to $z-1$ arcs. Hence
\begin{equation}
P_1=\det(I-Y'),
\label{eq:md-Y}
\end{equation}
with $Y'$ equal to $Y$ minus one arc. An internal consistency check is that $P_1$ must not depend on which
bond is kept. The check is nontrivial when the bonds at a site are inequivalent, and it holds for every result below.

\emph{Corollary (UST leaves).} Since $P_1$ is independent of the kept bond, the probability that a
degree-$z$ vertex is a leaf of the uniform spanning tree is exactly $zP_1$. On the square lattice this
reproduces the known value $(8/\pi^2)(1-2/\pi)$.

\subsection{Watson integrals}

Table~\ref{tab:watson} lists the three constants we need. For FCC we use
$W_f=\E[1/(3-s)]$, which is one third of Watson's normalization $\E[1/(1-s/3)]$.

\begin{table*}
\caption{Watson integrals ($c_i=\cos k_i$, $s=c_1c_2+c_2c_3+c_3c_1$).}
\label{tab:watson}
\begin{ruledtabular}
\begin{tabular}{lll}
Lattice & Definition and closed form & Value\\
\hline
SC & $\E\bigl[\tfrac{1}{1-(c_1+c_2+c_3)/3}\bigr]=\tfrac{\sqrt6\,\Gamma(\frac1{24})\Gamma(\frac5{24})\Gamma(\frac7{24})\Gamma(\frac{11}{24})}{32\pi^3}$ & $1.51638605915$\\
BCC & $\E\bigl[\tfrac{1}{1-c_1c_2c_3}\bigr]=\tfrac{\Gamma(\frac14)^4}{4\pi^3}$ & $1.39320392969$\\
FCC & $\E\bigl[\tfrac{1}{3-s}\bigr]=\tfrac{3\,\Gamma(\frac13)^6}{2^{14/3}\pi^4}$ & $0.44822039439$\\
\end{tabular}
\end{ruledtabular}
\end{table*}

\subsection{Resistances and results}

In each case below, $q\equiv1/(\pi^2W)$ with $W$ the lattice's own Watson integral.

\emph{Simple cubic} ($z=6$). The resistances are $R(1,0,0)=1/3$, $R(1,1,0)=\alpha$ and
$R(2,0,0)=2-4\alpha$, with
\begin{equation}
\alpha=\frac{7W+54/(\pi^2W)}{36}=0.395079152341851372\ldots
\end{equation}
The $5\times5$ determinant factorizes:
\begin{equation}
P_1^{\rm SC}=\tfrac43\,\alpha^3(2-3\alpha)^2=0.054582470221053656104\ldots
\end{equation}

\emph{Body-centred cubic} ($z=8$, neighbours $(\pm1,\pm1,\pm1)$). The resistances are
$R(1,1,1)=1/4$, $R(2,0,0)=(W-4q)/4$, $R(2,2,0)=4q$ and $R(2,2,2)=(8-3W-36q)/4$. Then
\begin{align}
P_1^{\rm BCC}&=\frac{3q}{128}(W+4q)^3(4-W-12q)^3\nonumber\\
&=0.042454248763323867910\ldots
\end{align}

\emph{Face-centred cubic} ($z=12$, neighbours are the permutations of $(\pm1,\pm1,0)$). The resistances are
$R(1,1,0)=1/6$, $R(2,0,0)=(4W-3q)/6$, $R(2,1,1)=(2W+3q-1)/3$ and $R(2,2,0)=(8-12W-9q)/3$. Then
\begin{align}
P_1^{\rm FCC}&=6^{-10}(5-4W-6q)^2(7-8W-3q)^3\nonumber\\
&\quad\times(8W+3q+1)^3(16W+15q-5)^3\nonumber\\
&=0.029126520511690534744\ldots
\end{align}

\emph{Pyrochlore} ($z=6$, corner-sharing tetrahedra, $W=W_f$, $X=1/(\pi^2W_f)$). The Laplacian
transfer currents from an edge are $1/3$ (self), $1/6$ (the two other edges of the same
tetrahedron), $(1+8W+3X)/48$ (collinear partner) and $(15-8W-3X)/96$ (the other two). Then
Eq.~\eqref{eq:md-Y} gives a perfect square:
\begin{align}
P_1^{\rm pyro}&=\Bigl[\frac{(31-8W-3X)(67+24W+9X)}{9216}\Bigr]^2\nonumber\\
&=0.053582950138955615678\ldots
\end{align}

\emph{Diamond} follows from the general theorem of Sec.~\ref{sec:bip}: $P_1=2/27$.

\subsection{Independent verification}

We computed $P_1$ \emph{without sampling} at the centre of finite boxes with dissipative
boundaries, by exact linear algebra (no sampling). The sequences converge to the formulas:
\begin{center}
\footnotesize
\begin{tabular}{lll}
\toprule
Lattice & Boxes & Exact finite-box $P_1$\\
\midrule
SC & $N=16,24,32$ & $0.0546206,\ 0.0545941,\ 0.0545875$\\
BCC & $N=16,24,32$ & $0.0424902,\ 0.0424653,\ 0.0424590$\\
FCC & $N=12,18,24$ & $0.0291795,\ 0.0291427,\ 0.0291335$\\
FCC & $N=30$ & $0.0291301$\\
Pyrochl. & $N=10,16,22$ & $0.0536309,\ 0.0535949,\ 0.0535876$\\
Diamond & $N$ increasing & $0.0742235,\ 0.0741125,\ 0.0740859$\\
Laves & $N=10,16,24$ & $0.0834488,\ 0.0833615,\ 0.0833418$\\
\bottomrule
\end{tabular}
\end{center}
Monte Carlo simulations on $32^3$--$44^3$ lattices give SC $0.05451(23)$, BCC $0.04244(21)$ and
FCC $0.02926(8)$. An earlier FCC run gave $0.02865(17)$, which we attribute to slow mixing; the exact
finite-box sequence resolves the discrepancy in favour of the formula. The corresponding UST leaf
probabilities $zP_1$ are SC $0.3274948$, BCC $0.3396340$, FCC $0.3495182$, pyrochlore $0.3214977$
and diamond $8/27$.

%=====================================================================
\section{Exact $P(h=1)$ calculations}
\label{sec:p1}
%=====================================================================

\subsection{A general theorem for bipartite lattices}
\label{sec:bip}

\emph{Theorem 1.} Let the lattice be $z$-regular, bipartite, of girth at least 6 (so that
$A^2=zI+A_2$), and edge-transitive, and let its second-neighbour graph $A_2$ also be edge-transitive.
Then $R(\text{nearest})=2/z$, $R(\text{second})=2/(z-1)$, and
\begin{equation}
P_1=\frac{(z-2)^{z-1}}{z\,(z-1)^{z-1}} .
\label{eq:bip}
\end{equation}
\emph{Proof.} Write $L_2=z^2I-A^2=LQ$. This operator acts within each sublattice, so
$R_L(u,v)=z\,R_{L_2}(u,v)$ for $u,v$ on the same sublattice. Foster's theorem~\cite{Foster1949} on
the edge-transitive second-neighbour graph gives $R_{L_2}=2/[z(z-1)]$ across a second-neighbour bond.
Hence $Y'$ has the constant diagonal $2/z$ and constant off-diagonal entries. Its determinant then
factorizes into Eq.~\eqref{eq:bip}. $\square$

The theorem gives $1/12$ for the honeycomb, which is known~\cite{AzimiTafreshi2010}, and $1/12$ for
the Laves graph (srs). It gives $2/27=0.0740740\ldots$ for diamond and applies to $d$-dimensional
diamond with $z=d+1$.

\subsection{The king lattice}

Take $\mathbb{Z}^2$ with all eight king-move neighbours ($z=8$). Put $p=1/\pi$ and
$L=\ln(2+\sq)/(\pi\sq)=0.24202564754285515484$. The resistances are
\begin{align}
R(1,0)&=L, & R(1,1)&=\tfrac12-L,\nonumber\\
R(2,0)&=\tfrac4\pi-4L, & R(2,1)&=6L-\tfrac12-\tfrac2\pi,\nonumber\\
R(2,2)&=5-\tfrac4\pi-14L,
\end{align}
and the $7\times7$ determinant factorizes as
\begin{align}
P_1^{\rm king}&=-\tfrac1{2048}(3+4p-14L)(1+4p-2L)(6L+12p-7)\nonumber\\
&\quad\times(13-32p+48p^2-52L-12L^2)^2\nonumber\\
&=0.0420560364703691189817\ldots
\end{align}
A Monte Carlo simulation on $96\times96$ gives $0.04218\pm0.00033$.

\subsection{Further two-dimensional lattices}

Table~\ref{tab:p1} collects the remaining exact values. We write
$c=\sqrt2\arccos(1/3)/\pi=0.55412642397957198712$ for the 4.8.8/Union Jack constant, $p=1/\pi$ and
$r=\sq/\pi$.

\begin{table*}
\caption{Exact height-one probabilities $P_1$ of the Abelian sandpile. Each value is independent
of which bond is kept in Eq.~\eqref{eq:md}.}
\label{tab:p1}
\begin{ruledtabular}
\begin{tabular}{llr}
Lattice (site, degree) & $P_1$ & Value\\
\hline
4.8.8 (3) & $c(4-3c)/16$ & $0.08095858841632017693$\\
Union Jack, face centre (4) & $9c^2(1-c)/8$ & $0.15402172338027135905$\\
Union Jack, corner (8) & $\tfrac1{256}(9c-2-4p)(6-7c-4p)(7-9c-4p)(81c^2-105c+34-4p-16p^2)^2$ & $0.02043343799324922073$\\
Star 3.12.12 (3) & $7/100$ & $0.07$\\
Dice, rim (3) & $3/16$ & $0.1875$\\
Dice, hub (6) & $\bigl[(1+2r)(11-6r)/144\bigr]^2$ & $0.01261518916334840149$\\
Lieb, corner (4) & $(\pi-2)/(4\pi^3)$ & $0.00920452869398469827$\\
Lieb, edge (2) & $1/4$ & $0.25$\\
Snub square (5) & $\tfrac{8953}{3888}-\tfrac{281\sq}{216}+\bigl(\tfrac{44\sq}{27}-\tfrac{11}4\bigr)p+\bigl(\tfrac{17}{36}-\tfrac{\sq}3\bigr)p^2$ & $0.06191455535342906165$\\
Checkerboard (6) & $13(7-8p)(11+8p)(1+16p-16p^2)/65536$ & $0.05351525370368361765$\\
Triakis, lattice vertex (12) & $(6r-1)(5+18r)^2(41-42r)^2/45562500$ & $0.00359238258113540057$\\
Triakis, triangle centre (3) & $27/100$ & $0.27$\\
Kagome (4) [known~\cite{Liu2026}] & $(5+3r)(17-6r)/1296$ & $0.07029827383269397637$\\
Laves graph (3) & $1/12$ & $0.08\overline{3}$\\
Diamond (4) & $2/27$ & $0.\overline{074}$\\
\end{tabular}
\end{ruledtabular}
\end{table*}

\emph{Honeycomb with next-nearest neighbours} (unit weights, $z=9$). Let
$g=4K(64/189)/(3\sqrt{21}\pi)=0.16088977460548913890$, the triangular Green function at $t=13/2$,
and let $g'=dg/dt$. All resistances among the site and its nine neighbours are rational combinations
of $1,r,g,g'$, and
\begin{align}
P_1&=\frac{(2-7g)(12r-14-105g-1120g')\,\mathcal P^2}{196036848},\nonumber\\
\mathcal P&=729120gg'+1764gr^2-7224gr+9163g+8202600g'^2\nonumber\\
&\quad-136080g'r+305760g'+792r^2-2688r+2548,
\end{align}
which evaluates to $P_1=0.037899503465150422174$. Exact finite boxes of $40^2$, $80^2$ and $160^2$ cells give
$0.0379335$, $0.0379082$ and $0.0379017$.

\emph{Validation against the literature.} The same pipeline gives the kagome value
$(5+3r)(17-6r)/1296$, which matches the value reported in Ref.~\cite{Liu2026}. It gives $1/12$ for the
honeycomb~\cite{AzimiTafreshi2010} and reproduces the square lattice's $2/\pi^2-4/\pi^3$ to 20 digits.
Monte Carlo checks on open patches: 4.8.8 $0.0807$--$0.0817\,(\pm0.0006)$; star
$0.0697$--$0.0706\,(\pm0.0008)$; dice hub $0.01273(30)$, rim $0.1888(10)$; Union Jack corner
$0.02024(30)$, centre $0.1535(8)$. The checkerboard and triakis values were confirmed by exact
finite boxes ($40^2,80^2,160^2$ cells): $0.0535660$, $0.0535281$, $0.0535185$ and
$0.0035988$, $0.0035940$, $0.0035928$.

%=====================================================================
\section{Spanning-tree entropy on Cairo and snub-square lattices}
\label{sec:entropy}
%=====================================================================

The number of spanning trees of an $N$-vertex piece of a lattice grows as $e^{zN}$~\cite{Wu1977}. For a
periodic lattice, $z$ is $|V_{\rm cell}|^{-1}\E_k\ln\det L(k)$.

\subsection{Snub square and Cairo}

For the snub square tiling $3^2.4.3.4$ (the Shastry--Sutherland lattice, 4 vertices per cell),
\begin{align}
\det L(k)&=2\bigl(170-72(\cos k_1+\cos k_2)-28\cos k_1\cos k_2\nonumber\\
&\qquad+\cos2k_1+\cos2k_2\bigr).
\end{align}
We find
\begin{equation}
\E_k\ln\det L=\frac{40}{3\pi}G+\frac43\ln(2+\sq),
\end{equation}
and hence
\begin{align}
z_{\rm snub}&=\frac{10}{3\pi}G+\frac13\ln(2+\sq)\nonumber\\
&=1.41085564574433483667\ldots,\\
z_{\rm Cairo}&=\frac{20}{9\pi}G+\frac29\ln(2+\sq)\nonumber\\
&=0.94057043049622322445\ldots
\end{align}
The Cairo pentagonal tiling is the planar dual. It has the same number of trees per cell but 6 vertices per cell.

\emph{Why it closes.} The $z_2$-discriminant of the spectral curve $\det L=0$ has only one simple factor,
$z_1^2-98z_1+1$, whose roots are $(\sq\pm\sqrt2)^4$. After the node at $k=0$ is removed, the curve therefore has genus zero. The entropy must then be
a combination of dilogarithms at algebraic points. PSLQ over
$\{\ln2,\ln3,\ln(1+\sqrt2),\ln(2+\sq),\ln(\sqrt2+\sq),G/\pi,\Cl(\pi/3)/\pi,\Cl(\pi/12)/\pi,\Cl(5\pi/12)/\pi\}$
returned $-3\,\E\ln\det L+4\ln(2+\sq)+40G/\pi=0$~\cite{PSLQ1999} with residual below $10^{-58}$ at 60 digits. The numerical value
$1.41086$ is known~\cite{ChangWang2006}; we found no prior closed form.

The same genus-zero structure gives \emph{algebraic} resistances (Table~\ref{tab:res}):
$R_{\rm tri\text{-}sq}=\tfrac12-\sq/18$ and $R_{\rm tri\text{-}tri}=2\sq/9$. Foster's sum rule
$8(\tfrac12-\sq/18)+2(2\sq/9)=4$ holds exactly. Planar duality, $R(e)+R(e^*)=1$ (an edge is in the
spanning tree exactly when its dual edge is not in the dual tree), then gives the Cairo values
$\tfrac12+\sq/18$ and $1-2\sq/9$. The remaining transfer currents at a snub-square vertex are
$Y(a,b)=1-\sq/2$, $Y(a,d)=\sq/9$, $Y(a',b')=(3-\sq)/6$, $Y(d,a')=\tfrac12-2\sq/9$,
$Y(a,a')=5\sq/9-1+1/(2\pi)$ and $Y(a,b')=\tfrac12-\sq/9-1/(2\pi)$. Here $d$ is the dimer arc; $a,b$ each
form a triangle with $d$; $a',b'$ form a triangle with each other; $a'$ is collinear with $a$ and $b'$ with $b$.

\subsection{Shastry--Sutherland network with arbitrary dimer conductance}

Give the square bonds conductance 1 and the dimer bonds conductance $w\ge0$, and let
$\Phi=\arccos\bigl(1/(w+1)\bigr)$. Then
\begin{align}
R_{\rm dimer}&=\frac{2\Phi}{\pi\sqrt{w(w+2)}},\\
R_{\rm square}&=\frac12-\frac{\sqrt w\,\Phi}{2\pi\sqrt{w+2}},\\
z(w)&=\frac1\pi\Bigl[2G+\Phi\arccosh(w+1)\nonumber\\
&\qquad+\Cl\bigl(\tfrac\pi2+\Phi\bigr)+\Cl\bigl(\tfrac\pi2-\Phi\bigr)\Bigr].
\label{eq:ssz}
\end{align}
\emph{Derivation.} Start from $z(0)=4G/\pi$. Kirchhoff's relation
$\partial z/\partial w=\tfrac12R_{\rm dimer}(w)$ (the derivative of the log tree count with respect to a
conductance is the sum of the corresponding edge resistances) then integrates via $u+1=\sec\phi$ and
\begin{align}
\int_0^\Phi\phi\sec\phi\,d\phi&=\Phi\ln\tan\bigl(\tfrac\pi4+\tfrac\Phi2\bigr)-2G\nonumber\\
&\quad+\Cl\bigl(\tfrac\pi2+\Phi\bigr)+\Cl\bigl(\tfrac\pi2-\Phi\bigr).
\end{align}
\emph{Special cases.} At $w=1$, $\Phi=\pi/3$, and $\Cl(\pi/6)+\Cl(5\pi/6)=4G/3$ recovers the snub-square
values. At $w=2$, $z(2)=1.57336855075766435824$, which coincides with the independently computed
Union Jack entropy per vertex. As $w\to\infty$, $z=\tfrac12\ln(2w)+2G/\pi+o(1)$.

\subsection{The king lattice}

Take $\mathbb{Z}^2$ with unit axis bonds and diagonal conductance $w$, and let $\varphi=\arccos(2w)$. Then
\begin{align}
R_{\rm axis}&=\frac{\arccos(2w)}{\pi\sqrt{1-4w^2}},\\
R_{\rm diag}&=\frac1{2w}-\frac{\arccos(2w)}{\pi w\sqrt{1-4w^2}},
\end{align}
for $w<\tfrac12$. For $w>\tfrac12$, use the continuation $\arccosh(2w)/\sqrt{4w^2-1}$. The entropy is
\begin{align}
z(w)&=\frac{4G}\pi+\ln w+\frac2\pi\Bigl[\varphi\ln\tan\bigl(\tfrac\pi4+\tfrac\varphi2\bigr)\nonumber\\
&\qquad+\Cl\bigl(\tfrac\pi2+\varphi\bigr)+\Cl\bigl(\tfrac\pi2-\varphi\bigr)\Bigr],
\end{align}
and at $w=1$ (the plain king graph)
\begin{equation}
z_{\rm king}=\frac{20G}{3\pi}=1.94373936020545853426\ldots,
\end{equation}
with $R_{\rm axis}=\ln(2+\sq)/(\pi\sq)$ and $R_{\rm diag}=\tfrac12-R_{\rm axis}$. The determinant
$9-(1+x+x^{-1})(1+y+y^{-1})$ belongs to a family of Mahler measures, so $20G/(3\pi)$ may be implicit
in that literature~\cite{GuttmannRogers2012}. The elementary structure has a simple origin: after
the exact $k_2$ average and the substitution $x=-1+t^2$, the remaining integrand has discriminant
$8w+(2w-1)^2=(2w+1)^2$, a perfect square.

\emph{Corollaries.} (i) Eliminating each face centre of the Union Jack lattice with centre-bond
conductance $w$ (star--mesh) yields the king lattice with axis conductance $1+w/2$ and diagonal
ratio $w'=w/(4+2w)$. At $w=1$ this reproduces $R_{\rm corner\text{-}corner}=c/2$.
(ii) The checkerboard lattice with crossed-diagonal conductance $w$ has
$\det=4(w+1)[2(a+b)+(w-1)ab]$, with $a=1-\cos(k_1+k_2)$ and $b=1-\cos(k_1-k_2)$. This is the king
determinant with $w'=(1-w)/(2(1+w))$.

%=====================================================================
\section{Resistor network reductions}
\label{sec:res}
%=====================================================================

\subsection{Adjacent resistances on Archimedean and Laves tilings}

\begin{table}
\caption{Exact adjacent-node resistances ($1\,\Omega$ resistors). $c=\sqrt2\arccos(1/3)/\pi$.
Pairs related by planar duality satisfy $R(e)+R(e^*)=1$.}
\label{tab:res}
\begin{ruledtabular}
\begin{tabular}{lll}
Tiling, edge & $R$ & Value\\
\hline
Snub square, tri--sq & $\tfrac12-\tfrac{\sq}{18}$ & $0.40377495513506237$\\
Snub square, tri--tri & $\tfrac{2\sq}{9}$ & $0.38490017945975051$\\
Cairo, dual of tri--sq & $\tfrac12+\tfrac{\sq}{18}$ & $0.59622504486493763$\\
Cairo, dual of tri--tri & $1-\tfrac{2\sq}{9}$ & $0.61509982054024949$\\
4.8.8, square & $\tfrac12+\tfrac c4$ & $0.63853160599489299$\\
4.8.8, oct--oct & $1-\tfrac c2$ & $0.72293678801021400$\\
Union Jack, corner--corner & $\tfrac c2$ & $0.27706321198978599$\\
Union Jack, corner--centre & $\tfrac{2-c}{4}$ & $0.36146839400510700$\\
King, axis & $\tfrac{\ln(2+\sq)}{\pi\sq}$ & $0.24202564754285515$\\
King, diagonal & $\tfrac12-\tfrac{\ln(2+\sq)}{\pi\sq}$ & $0.25797435245714485$\\
\end{tabular}
\end{ruledtabular}
\end{table}

The 4.8.8 and Union Jack lattices are dual, and the entries in Table~\ref{tab:res} satisfy
$R(e)+R(e^*)=1$ identically. The 4.8.8 determinant factorizes as $64-4(3+\cos k_1)(3+\cos k_2)$,
which is the origin of $\arccos(1/3)$. Owaidat computed Union Jack resistances numerically~\cite{Owaidat2013};
Cserti \emph{et al.}\ gave exact mappings for kagome, dice and decorated lattices~\cite{Cserti2011}.
Kirchhoff solves on $120\times120$-cell tori, Richardson-extrapolated, agree with the table to about
$10^{-7}$.

\subsection{Square lattice with arbitrary diagonal bonds}

Give conductance $w_1$ to bonds along $e_1$, $w_2$ along $e_2$, $w_+$ along $e_1+e_2$ and $w_-$ along
$e_1-e_2$. Let $\sigma=w_++w_-$ and $\Delta=4w_+w_--w_1^2$. Then
\begin{equation}
R(e_1)=\begin{cases}
\dfrac{2}{\pi\sqrt\Delta}\arcsinh\sqrt{\dfrac{\Delta}{(w_1+\sigma)(w_1+w_2)}}, & \Delta>0,\\[2ex]
\dfrac{2}{\pi\sqrt{-\Delta}}\arcsin\sqrt{\dfrac{-\Delta}{(w_1+\sigma)(w_1+w_2)}}, & \Delta<0,\\[2ex]
\dfrac{2}{\pi\sqrt{(w_1+\sigma)(w_1+w_2)}}, & \Delta=0 .
\end{cases}
\end{equation}
\emph{Derivation.} At fixed $x=\cos k_1$ the $k_2$ average of $1/\det$ is $1/\sqrt{A^2-M^2}$ with
$A^2-M^2=4(1-x)(p+qx)$, where $p=(w_1+w_2+\sigma)^2-w_2^2-(w_+-w_-)^2$ and $q=\Delta$. The factor $1-x$ is
the node at $k=0$. Then $R(e_1)=\pi^{-1}\int_{-1}^1dx/\sqrt{(1+x)(p+qx)}$, and $p-q=2(w_1+\sigma)(w_1+w_2)$.
\emph{Special cases.} The square lattice gives $1/2$; the king lattice gives $\arccos(2w)$ as above; and the
triangular lattice ($w_-=0$) gives Kenyon's isoradial value~\cite{Kenyon2002}, $1/3$ in the isotropic case.
Together with Foster's identity $w_1R(e_1)+w_2R(e_2)+w_+R(e_1+e_2)+w_-R(e_1-e_2)=1$, this also gives the
diagonal resistances when $w_+=w_-$. The formula agrees with direct 2D integration to $\sim10^{-26}$
for random conductances.

\subsection{Next-nearest-neighbour networks on bipartite lattices}

Let the lattice be $z$-regular, bipartite and of girth at least 6, with unit nearest-neighbour
bonds and conductance $w$ on second neighbours. Since $A^2=zI+A_2$,
\begin{equation}
L=(z+z(z-1)w)I-A-w(A^2-zI),
\end{equation}
\begin{equation}
\det=u(1+2zw+w^2u),\qquad u=z^2-|f(k)|^2 .
\end{equation}
The symmetry identity $\E[\mathrm{Re}f/\det]=\E[|f|^2/(z\det)]$ collapses the nearest-neighbour resistance to
\begin{align}
R_{\rm NN}(w)&=2\Bigl(w+\frac1z\Bigr)\E_k\Bigl[\frac1{1+2zw+w^2(z^2-|f(k)|^2)}\Bigr],\\
R_{\rm NNN}(w)&=\frac{2-zR_{\rm NN}}{z(z-1)w}.
\end{align}
\emph{Honeycomb} ($z=3$, $|f|^2=3+2S$, $S=\cos k_1+\cos k_2+\cos(k_1-k_2)$). Here
$\det=2T(6w+1+2w^2T)$ with $T=3-S$. Horiguchi's formula~\cite{Horiguchi1972} with $e=3+1/w$ gives
\begin{align}
m&=\frac{16w^3(3w+1)}{(2w+1)^3(6w+1)},\\
R_{\rm NN}&=\frac{4(3w+1)K(m)}{3\pi(2w+1)^{3/2}\sqrt{6w+1}},
\end{align}
and $R_{\rm NNN}=(2-3R_{\rm NN})/(6w)$. At $w=1$: $m=64/189$,
$R_{\rm NN}=16K(64/189)/(9\sqrt{21}\pi)=0.21451969947398551853$ and
$R_{\rm NNN}=0.22607348359634057407$.
\emph{Diamond} ($z=4$). Here
\begin{equation}
R_{\rm NN}=\frac{4w+1}{8w^2}\,P_{\rm FCC}\Bigl(3+\frac{8w+1}{4w^2}\Bigr),
\end{equation}
where $P_{\rm FCC}(t)=\E[1/(t-s)]$ is the FCC Green function in Joyce's closed form~\cite{Joyce1998}. At $w=1$
($t=21/4$), $R_{\rm NN}=0.12347073437196593577$ and $R_{\rm NNN}=(1-2R_{\rm NN})/(6w)=0.12550975520934468$.
Three-dimensional Kirchhoff solves on $12^3$, $16^3$ and $20^3$ tori give
$0.1234547\to0.1234640\to0.1234673$.

%=====================================================================
\section{Discussion and conclusion}
\label{sec:conclusion}
%=====================================================================

Three observations organize the results. First, the Majumdar--Dhar matrix \emph{is} the
transfer-current matrix, Eq.~\eqref{eq:md-Y}. Sandpile height probabilities, UST leaf and degree
statistics, and resistances therefore all come from the same $z\times z$ matrix at a vertex. Second, in
three dimensions every needed resistance lies in the three-dimensional space spanned by
$\{1,W,1/(\pi^2W)\}$. This is why $P_1$ becomes a polynomial in the Watson integral, and why the
determinants factor so strikingly. Third, in two dimensions, closed forms arise whenever the spectral curve has
genus zero after removal of the node at $k=0$. That gives algebraic or elementary resistances and
Clausen-type entropies (snub square, Cairo, Shastry--Sutherland, king). Otherwise complete elliptic
integrals appear (next-nearest-neighbour honeycomb).

Open questions include the higher height probabilities $P_2,P_3,\ldots$ in three dimensions, the
ruby and snub-hexagonal lattices (for which PSLQ found no closed form in natural bases), and a
conceptual explanation of the perfect-square factorization of the pyrochlore $P_1$.

\appendix

\section{Added mass of a sphere touching a spherical container}
\label{app:added}

A companion image-sum result is recorded here because it has the same flavour: an infinite
image series collapses to a single special value. Consider a ball of radius $a$ that touches, from
inside, a fixed spherical tank of radius $b>a$ filled with ideal (inviscid, incompressible,
irrotational) fluid. The ball moves along the line of centres, and its added mass is
$M=C\rho\tfrac43\pi a^3$. With $r=a/b$ and $q=1/(1-r)$,
\begin{equation}
C(r)=\frac32\sum_{k=0}^\infty\frac1{[1+k(1-r)]^3}-1=\frac32q^3\zeta(3,q)-1,
\end{equation}
where $\zeta(s,q)$ is the Hurwitz zeta function. More generally, let the ball touch a sphere of
\emph{signed} radius $b$ ($b>0$ a container, $b=\infty$ a plane, $b<0$ an external ball of radius
$|b|$) and set $\lambda=1-a/b$. Then $C=\tfrac32\sum_{k\ge0}(1+k\lambda)^{-3}-1$. Special cases:
a flat wall gives $\tfrac32\zeta(3)-1=0.8030853\ldots$ (known); $r=\tfrac12$ gives $12(\zeta(3)-1)-1=1.4246828\ldots$;
and two equal touching balls give $\tfrac{21}{16}\zeta(3)-1=0.5776997\ldots$

\emph{Derivation.} The image of an axial dipole in a rigid sphere is an axial dipole at the
inverse point with strength scaled by $-(R/f)^3$~\cite{Hicks1880,Lamb1932}. Inverting about the contact
point maps both spheres to parallel planes, so the image positions form an arithmetic progression.
The strength ratios then telescope:
$A_k=A_0\,b^3/[a+k(b-a)]^3$. The added mass, $-(\rho/U)\oint\phi\,n_x\,dS$ over the ball, is a sum of
Hurwitz-zeta values.

\emph{Verification.} We reconstructed the flow from the images and checked both boundary conditions at
2000 random surface points. Brute-force surface integration, with no zeta functions used, gives, for example,
$1.424682623$ ($N=6000$ images) against $1.424682838$ for $a=1,b=2$. The truncation gap shrinks by a factor
of about 4 per doubling of $N$. This result is probably a routine extension of classical image methods.

\begin{acknowledgments}
The computations and the first draft of this manuscript were carried out with the assistance
of an AI coding agent (Claude Code). All numerical code is available in the accompanying
repository.
\end{acknowledgments}

\begin{thebibliography}{99}

\bibitem{BTW1987}
P.~Bak, C.~Tang, and K.~Wiesenfeld,
Self-organized criticality: An explanation of the $1/f$ noise,
\href{https://doi.org/10.1103/PhysRevLett.59.381}{Phys.\ Rev.\ Lett.\ \textbf{59}, 381 (1987)}.

\bibitem{Dhar1990}
D.~Dhar,
Self-organized critical state of sandpile automaton models,
\href{https://doi.org/10.1103/PhysRevLett.64.1613}{Phys.\ Rev.\ Lett.\ \textbf{64}, 1613 (1990)};
Erratum, Phys.\ Rev.\ Lett.\ \textbf{64}, 2837 (1990).

\bibitem{MajumdarDhar1991}
S.~N.~Majumdar and D.~Dhar,
Height correlations in the Abelian sandpile model,
\href{https://doi.org/10.1088/0305-4470/24/7/008}{J.\ Phys.\ A: Math.\ Gen.\ \textbf{24}, L357 (1991)}.

\bibitem{Priezzhev1994}
V.~B.~Priezzhev,
Structure of two-dimensional sandpile. I. Height probabilities,
\href{https://doi.org/10.1007/BF02188212}{J.\ Stat.\ Phys.\ \textbf{74}, 955 (1994)}.

\bibitem{AzimiTafreshi2010}
N.~Azimi-Tafreshi, H.~Dashti-Naserabadi, S.~Moghimi-Araghi, and P.~Ruelle,
The Abelian sandpile model on the honeycomb lattice,
\href{https://doi.org/10.1088/1742-5468/2010/02/P02004}{J.\ Stat.\ Mech.\ (2010) P02004},
arXiv:0912.3331.

\bibitem{PonceletRuelle2017}
A.~Poncelet and P.~Ruelle,
Sandpile probabilities on triangular and hexagonal lattices,
\href{https://doi.org/10.1088/1751-8121/aa9255}{J.\ Phys.\ A: Math.\ Theor.\ \textbf{51}, 015002 (2018)},
arXiv:1708.03493.

\bibitem{Liu2026}
Q.~Liu, J.-N.~Herre, P.~Baruah, S.~Dreizler, C.~Karrasch, W.~Ruszel, and D.~Schuricht,
Universal correlations in the Abelian sandpile model,
arXiv:2609.10352 (2026).

\bibitem{Watson1939}
G.~N.~Watson,
Three triple integrals,
\href{https://doi.org/10.1093/qmath/os-10.1.266}{Q.\ J.\ Math.\ \textbf{os-10}, 266 (1939)}.

\bibitem{GlasserZucker1977}
M.~L.~Glasser and I.~J.~Zucker,
Extended Watson integrals for the cubic lattices,
\href{https://doi.org/10.1073/pnas.74.5.1800}{Proc.\ Natl.\ Acad.\ Sci.\ USA \textbf{74}, 1800 (1977)}.

\bibitem{Kirchhoff1847}
G.~Kirchhoff,
\"Uber die Aufl\"osung der Gleichungen, auf welche man bei der Untersuchung der linearen
Vertheilung galvanischer Str\"ome gef\"uhrt wird,
\href{https://doi.org/10.1002/andp.18471481202}{Ann.\ Phys.\ Chem.\ \textbf{72}, 497 (1847)}.

\bibitem{BurtonPemantle1993}
R.~Burton and R.~Pemantle,
Local characteristics, entropy and limit theorems for spanning trees and domino tilings via
transfer-impedances,
\href{https://doi.org/10.1214/aop/1176989121}{Ann.\ Probab.\ \textbf{21}, 1329 (1993)}.

\bibitem{LyonsPeres2016}
R.~Lyons and Y.~Peres,
\emph{Probability on Trees and Networks},
Cambridge Series in Statistical and Probabilistic Mathematics Vol.~42
(Cambridge University Press, Cambridge, 2016).

\bibitem{ChangWang2006}
S.-C.~Chang and W.~Wang,
Spanning trees on lattices and integration identities,
J.\ Phys.\ A: Math.\ Gen.\ \textbf{39}, 10263 (2006),
arXiv:cond-mat/0604589.

\bibitem{ChangShrock2006}
S.-C.~Chang and R.~Shrock,
Some exact results for spanning trees on lattices,
J.\ Phys.\ A: Math.\ Gen.\ \textbf{39}, 5653 (2006),
arXiv:cond-mat/0602574.

\bibitem{GuttmannRogers2012}
A.~J.~Guttmann and M.~D.~Rogers,
Spanning tree generating functions and Mahler measures,
\href{https://doi.org/10.1088/1751-8113/45/49/494001}{J.\ Phys.\ A: Math.\ Theor.\ \textbf{45}, 494001 (2012)}.

\bibitem{Owaidat2013}
M.~Owaidat,
On the uniform tiling with electrical resistors,
arXiv:1305.5839 (2013).

\bibitem{Cserti2011}
J.~Cserti, G.~Sz\'echenyi, and G.~D\'avid,
Uniform tiling with electrical resistors,
J.\ Phys.\ A: Math.\ Theor.\ \textbf{44}, 215201 (2011),
arXiv:1102.4537.

\bibitem{Foster1949}
R.~M.~Foster,
The average impedance of an electrical network,
in \emph{Reissner Anniversary Volume: Contributions to Applied Mechanics}
(J.~W.~Edwards, Ann Arbor, 1949), pp.~333--340.

\bibitem{Wu1977}
F.~Y.~Wu,
Number of spanning trees on a lattice,
J.\ Phys.\ A: Math.\ Gen.\ \textbf{10}, L113 (1977).

\bibitem{Kenyon2002}
R.~Kenyon,
The Laplacian and Dirac operators on critical planar graphs,
\href{https://doi.org/10.1007/s00222-002-0249-4}{Invent.\ Math.\ \textbf{150}, 409 (2002)}.

\bibitem{Horiguchi1972}
T.~Horiguchi,
Lattice Green's functions for the triangular and honeycomb lattices,
\href{https://doi.org/10.1063/1.1666155}{J.\ Math.\ Phys.\ \textbf{13}, 1411 (1972)}.

\bibitem{Joyce1998}
G.~S.~Joyce,
On the cubic modular transformation and the cubic lattice Green functions,
J.\ Phys.\ A: Math.\ Gen.\ \textbf{31}, 5105 (1998).

\bibitem{Hicks1880}
W.~M.~Hicks,
On the motion of two spheres in a fluid,
\href{https://doi.org/10.1098/rstl.1880.0013}{Philos.\ Trans.\ R.\ Soc.\ Lond.\ \textbf{171}, 455 (1880)}.

\bibitem{Lamb1932}
H.~Lamb,
\emph{Hydrodynamics}, 6th ed.\ (Cambridge University Press, Cambridge, 1932).

\bibitem{PSLQ1999}
H.~R.~P.~Ferguson, D.~H.~Bailey, and S.~Arno,
Analysis of PSLQ, an integer relation finding algorithm,
\href{https://doi.org/10.1090/S0025-5718-99-00995-3}{Math.\ Comp.\ \textbf{68}, 351 (1999)}.

\end{thebibliography}

\end{document}
